\begin{corollary}[5.6.6]
\label{IV.2.s5.6.6-fr}
Bajo las hipótesis generales de \textup{(5.6.5)}, se tiene
\begin{equation}
\dim(X)\leq\dim(Y)+e.
\tag{5.6.6.1}\label{IV.2.s5.6.6.1-fr}
\end{equation}
Si se supone además que $Y$ es universalmente catenario, para que los dos
miembros de \textup{(5.6.6.1)} sean iguales, es necesario y suficiente que se tenga
\begin{equation}
\dim(Y)=\sup_{y\in f(X)}\dim(\mathcal O_y).
\tag{5.6.6.2}\label{IV.2.s5.6.6.2-fr}
\end{equation}
En particular, los dos miembros de \textup{(5.6.6.1)} son iguales cuando $Y$ es
localmente de tipo finito sobre un cuerpo.
\end{corollary}

La desigualdad (5.6.6.1) resulta, en efecto, de (5.6.5.1) aplicado en un punto
cerrado $x$ de $X$, teniendo en cuenta (5.1.4.2) y (5.1.11). Por otra parte,
toda fibra no vacía $f^{-1}(y)$ contiene un punto $x$ cerrado en esta
fibra (5.1.11), y si se supone que $Y$ es universalmente catenario, se tiene
por tanto en este punto la relación (5.6.5.3). Tomando los supremos de los dos
miembros de (5.6.5.3) cuando $x$ recorre $X$, y observando que todo punto
$x$ cerrado en $X$ es \emph{a fortiori} cerrado en $f^{-1}(f(x))$, la
segunda afirmación del corolario resulta de (5.1.4.1) y (5.1.4.2).

Para probar la última afirmación, observemos que existe una vecindad abierta
afín $U$ del punto genérico de $X$ tal que $f|U$ es de tipo finito, por tanto
$f(U)$ es una parte \emph{construible} de $Y$, densa en $Y$ (1.8.5), y por
consiguiente contiene un abierto no vacío (por tanto denso) de $Y$
$(\mathbf{0}_{\mathrm{III}},9.2.2)$. La hipótesis de que $Y$ es localmente de tipo
finito sobre un cuerpo asegura entonces, por una parte, que $Y$ es
universalmente catenario (5.6.4), y por otra, que los dos miembros de (5.6.6.2)
son iguales (5.2.2 y 4.1.1.3).

Se observará que la condición (5.6.6.2) se verifica trivialmente cuando $f$
es \emph{sobreyectivo}.

\begin{corollary}[5.6.7]
\label{IV.2.s5.6.7-fr}
Sean $Y$ un preesquema localmente noetheriano, $f:X\to Y$ un morfismo
localmente de tipo finito y $n$ un entero. Si, para todo $y\in Y$,
$\dim(f^{-1}(y))\leq n$, entonces se tiene
\begin{equation}
\dim(X)\leq\dim(Y)+n.
\tag{5.6.7.1}\label{IV.2.s5.6.7.1-fr}
\end{equation}
\end{corollary}

En virtud de (\textbf 0, 14.1.7) y de (1.5.4), es posible limitarse al caso en que $X$ e
$Y$ son afines, por tanto $f$ es de tipo finito y $X$ e $Y$ son reducidos;
sean $X_i$ ($1\leq i\leq m$) los subpreesquemas cerrados (íntegros) de $X$
cuyos espacios subyacentes son las componentes irreducibles de $X$; se tiene
$\dim(X)=\sup_i(\dim(X_i))$. Si $Z_i$ es el subpreesquema cerrado reducido de
$Y$ cuyo espacio subyacente es $\overline{f(X_i)}$, $Z_i$ es íntegro
$(\mathbf{0}_{\mathrm I},2.1.5)$; la restricción $X_i\to Y$ de $f$ se factoriza en
$X_i\xrightarrow{g_i}Z_i\xrightarrow{j_i}Y$, donde $j_i$ es la inyección
canónica (\textbf I, 5.2.2), y $g_i$ es dominante
\oldpage[IV-2]{101}
y de tipo finito (1.5.4).

Se tiene $\dim(Z_i)\leq\dim(Y)$ para todo $i$; por otra parte, si $z_i$ es
el punto genérico de $Z_i$, se tiene $\dim(g_i^{-1}(z_i))\leq n$ para
todo $i$ por hipótesis; se ve así que para probar (5.6.7.1), basta demostrar
que $\dim(X_i)\leq\dim(Z_i)+n$ para todo $i$, lo que resulta de (5.6.6.1).

\begin{corollary}[5.6.8]
\label{IV.2.s5.6.8-fr}
Sean $Y$ un preesquema localmente noetheriano, $X$ un preesquema
\textup{irreducible}, $f:X\to Y$ un morfismo localmente de tipo finito y $n$ un entero
$\geq0$. Se supone que $Y$ es universalmente catenario y que, para todo
$y\in Y$, se tiene $\dim(f^{-1}(y))\geq n$ \textup{(resp.}
$\dim(f^{-1}(y))=n$\textup{)}. Entonces se tiene
\begin{equation}
\dim(X)\geq\dim(Y)+n
\tag{5.6.8.1}\label{IV.2.s5.6.8.1-fr}
\end{equation}
\textup{(resp.}
\begin{equation}
\dim(X)=\dim(Y)+n).
\tag{5.6.8.2}\label{IV.2.s5.6.8.2-fr}
\end{equation}
\end{corollary}

Como $n\geq0$, la hipótesis implica que $f$ es \emph{sobreyectivo}, por tanto $Y$
es irreducible; además, si $\eta$ es el punto genérico de $Y$, se tiene
$\dim(f^{-1}(\eta))\geq n$ (resp. $\dim(f^{-1}(\eta))=n$). La conclusión
resulta por tanto de (5.6.6).
